\documentclass[aspectratio=169]{beamer}
\input{header}
\coursecode{JKSSB}

\title{Virtual Memory}
\subtitle{Lesson 09 --- Paging, Segmentation and Swapping}
\author{JKSSB Computer Awareness}
\date{\today}

\begin{document}
\frame{\titlepage}

\begin{frame}{Why Virtual Memory Exists}
  \begin{itemize}
    \item A program or its data can be \key{larger than physical RAM}.
    \item Virtual memory lets the system run such programs anyway, using
      secondary storage (disk) as an extension.
    \item It gives each process a \key{large, contiguous view} of memory,
      even though the real RAM is smaller and shared.
    \item The office desktop (EX-01) runs MS Word, a spreadsheet, and the
      browser together; virtual memory keeps them all addressable.
  \end{itemize}
  \begin{block}{High-yield point}
    Virtual memory \key{extends} addressable memory using disk; it is
    \bad{not} extra physical RAM.
  \end{block}
\end{frame}

\begin{frame}{Virtual vs Physical Addresses}
  \begin{columns}[T]
    \column{0.46\textwidth}
      \textbf{Virtual address}
      \begin{itemize}
        \item The address a program generates.
        \item Large, uniform space per process.
        \item Not the real RAM location.
      \end{itemize}
    \column{0.46\textwidth}
      \textbf{Physical address}
      \begin{itemize}
        \item The actual location in RAM.
        \item Managed by the OS.
        \item Where the data really sits.
      \end{itemize}
  \end{columns}
  \vspace{0.2cm}
  \begin{block}{Exam cue}
    The program always sees \key{virtual} addresses; the hardware converts
    them to \key{physical} addresses.
  \end{block}
\end{frame}

\begin{frame}{The MMU: The Translator}
  \begin{itemize}
    \item \key{MMU} = \key{Memory Management Unit}.
    \item It is the hardware that \key{translates virtual addresses into
      physical addresses} at run time.
    \item It works with the page table the operating system maintains.
    \item Every memory access passes through this translation.
  \end{itemize}
  \begin{alertblock}{Trap}
    Translation is the MMU's job; \bad{storing files is not}. Virtual memory
    uses the disk, but the MMU only maps addresses.
  \end{alertblock}
\end{frame}

\begin{frame}{Paging: Fixed-Size Blocks}
  \begin{itemize}
    \item Paging divides the \key{virtual address space into fixed-size
      pages}.
    \item Physical memory is divided into equal \key{frames} of the same size.
    \item Any page can be loaded into any free frame.
    \item Because all blocks are equal, allocation is simple and fast.
  \end{itemize}
  \begin{block}{Exam cue}
    \key{Pages} are the virtual blocks; \key{frames} are the physical blocks.
    Both are fixed-size and equal.
  \end{block}
\end{frame}

\begin{frame}{Segmentation: Variable-Size Segments}
  \begin{itemize}
    \item Segmentation divides a program into \key{variable-size logical
      segments}.
    \item A segment matches a program unit: code, data, stack, a procedure,
      or a module.
    \item Segment sizes differ, because logical units differ.
    \item It reflects the \key{programmer's logical view}, not a fixed grid.
  \end{itemize}
  \begin{alertblock}{Trap}
    \bad{Segmentation is not paging.} Paging uses equal fixed-size pages and
    frames; segmentation uses unequal logical segments.
  \end{alertblock}
\end{frame}

\begin{frame}{Paging vs Segmentation}
  \begin{center}
    \begin{tabular}{ll}
      \toprule
      \textbf{Paging} & \textbf{Segmentation} \\
      \midrule
      Fixed-size blocks & Variable-size blocks \\
      Blocks: pages / frames & Blocks: logical segments \\
      Equal size & Unequal size \\
      Invisible to programmer & Matches program structure \\
      Simple allocation & Matches logical units \\
      \bottomrule
    \end{tabular}
  \end{center}
  \vspace{0.3cm}
  \muted{Remember: \textit{pages are fixed}; \textit{segments are logical and
  vary in size}.}
\end{frame}

\begin{frame}{Page Fault and Demand Paging}
  \begin{itemize}
    \item A \key{page fault} occurs when the referenced page is
      \key{not currently in RAM}.
    \item The page must then be \key{loaded from the disk} into a free frame.
    \item \key{Demand paging} loads a page only when it is actually needed,
      not all at once.
    \item This is why a program can start before all of it is in memory.
  \end{itemize}
  \begin{alertblock}{Trap}
    A page fault is \bad{not a hardware fault}. It is a normal signal that
    the needed page is currently on disk.
  \end{alertblock}
\end{frame}

\begin{frame}{Swapping and the Disk-Backed Hierarchy}
  \begin{itemize}
    \item \key{Swapping} moves pages or whole processes between \key{RAM and
      secondary storage (disk)}.
    \item Virtual memory is therefore \key{disk-backed}.
    \item The hierarchy from the CPU outwards: registers, cache, RAM, then
      disk-backed virtual memory.
    \item EX-01: with 16\,GB RAM, a working set that exceeds RAM spills to
      disk and is loaded back on demand.
  \end{itemize}
  \begin{block}{Exam cue}
    Disk-backed virtual memory is \key{much slower} than RAM. It adds
    \key{capacity}, not speed.
  \end{block}
\end{frame}

\begin{frame}{Not Extra RAM --- and Not Cache}
  \begin{itemize}
    \item Virtual memory is \bad{not additional physical RAM}; it borrows
      space on the disk.
    \item It is \bad{not extra speed}: reading from disk is far slower than
      reading from RAM.
    \item Do not confuse it with \key{cache}: cache is small, fast memory
      \key{close to the CPU}; virtual memory is \key{large and disk-backed}.
  \end{itemize}
  \begin{alertblock}{Common confusion}
    Cache \textit{speeds up} frequently used data; virtual memory
    \textit{extends capacity} using slower storage.
  \end{alertblock}
\end{frame}

\begin{frame}{Thrashing}
  \begin{itemize}
    \item \key{Thrashing} is \key{excessive paging}.
    \item The system spends most of its time moving pages between RAM and
      disk, leaving little time for useful work.
    \item It happens when too many processes compete for too little RAM.
    \item Performance collapses even though the CPU is busy.
  \end{itemize}
  \begin{block}{Exam cue}
    Thrashing = \key{too much paging}, not a hardware failure. Adding RAM
    usually relieves it.
  \end{block}
\end{frame}

\begin{frame}{Lesson 09: Recall}
  \begin{itemize}
    \item Virtual memory \key{extends memory using disk}; it is slower than
      RAM and is \key{not extra physical RAM}.
    \item \key{MMU} translates virtual addresses into physical addresses.
    \item \key{Paging} = fixed-size pages and frames; \key{segmentation} =
      variable-size logical segments.
    \item A \key{page fault} means the referenced page is not in RAM and must
      be loaded from disk; \key{demand paging} loads pages only when needed.
    \item \key{Thrashing} = excessive paging.
  \end{itemize}
\end{frame}
\end{document}
